% SPDX-License-Identifier: CC-BY-4.0
\documentclass[10.5pt]{article}

\usepackage[T1]{fontenc}
\usepackage[utf8]{inputenc}
\usepackage{lmodern}
\usepackage[margin=0.9in]{geometry}
\usepackage{amsmath,amssymb,amsthm}
\usepackage{booktabs}
\usepackage{microtype}
\usepackage{enumitem}
\usepackage{parskip}
\usepackage{array}
\usepackage{fancyvrb}
\usepackage{xcolor}
\usepackage[hidelinks]{hyperref}
\usepackage{xurl}
\urlstyle{same}

\hypersetup{
  pdftitle={Offline Evidence, Online Admission: Action-Bound Human Authorization Under a Signature-Soliciting Adversary},
  pdfauthor={Iman Schrock},
  pdfsubject={Cryptographic protocols and symbolic analysis},
  pdfkeywords={authorization, action binding, digital signatures, replay, Tamarin, AI agents}
}

\newtheorem{definition}{Definition}
\newtheorem{proposition}{Proposition}
\newtheorem{lemma}{Lemma}
\newcommand{\Enc}{\mathsf{Enc}}
\newcommand{\Hash}{\mathsf{H}}
\newcommand{\Sign}{\mathsf{Sign}}
\newcommand{\Verify}{\mathsf{Verify}}
\newcommand{\Accept}{\mathsf{Accept}}
\newcommand{\Consume}{\mathsf{Consume}}
\newcommand{\code}[1]{\texttt{#1}}

\title{Offline Evidence, Online Admission:\\
Action-Bound Human Authorization Under a Signature-Soliciting Adversary}

\author{Iman Schrock\\
EMILIA Protocol, Inc.\\
\texttt{team@emiliaprotocol.ai}}

\date{Preprint. August 2026.}

\begin{document}
\maketitle

\begin{abstract}
An agent can hold a valid session and broad delegated authority while still
needing a human decision for one consequential action discovered during
execution. A common response is to create a signed approval artifact that can be
checked later without contacting the approval service. Two different properties
are often attributed to that artifact: that it is bound to the action the human
approved, and that it can be used only once. This paper separates them.

We define an action-bound authorization receipt for an adversary that controls
the network and the approval orchestrator and may ask an honest approver to sign
arbitrary actions. Under standard signature and hash assumptions, offline
verification gives action agreement: a receipt accepted for action $a$ contains
an uncompromised approver's signature over a canonical commitment to $a$. It does
not give global one-time admission. We show that a static artifact presented to
two isolated, otherwise correct verifiers can be accepted once by each. Preventing
that replay requires a shared atomic consumption domain, equivalent coordination,
or a construction that restricts the artifact to one admission domain.

We encode both boundaries in Tamarin. The core model verifies
user-verification-gated action agreement and local injective acceptance, and finds the expected
replay when consumption is omitted. A second model verifies injectivity inside
each consumption domain while falsifying global injectivity across domains, even
without key compromise. The result is a narrow placement rule: signatures make
authorization evidence portable; state at the execution boundary makes admission
one-time.
\end{abstract}

\section{Introduction}

An AI agent may begin a task under a valid session and discover, only after it has
read data or planned a sequence of calls, that one step requires a human decision.
Examples include changing a payment destination, deleting a production resource,
or releasing confidential data. The question at that point is not simply whether
the agent has a credential. It is whether a particular human authorized a
particular action, and whether that authorization may still admit execution.

Digital signatures are a natural part of the answer. If the human signs a
canonical commitment to the action, a later verifier can test the signature and
the action binding without trusting the service that routed the approval. This is
useful when the operator of the approval service is also the party that would
otherwise produce the audit record. A portable signed artifact turns one claim
from an operator-controlled database entry into independently checkable evidence.

The same artifact is sometimes also described as ``single use'' or ``exactly
once.'' That language crosses a boundary. Signature verification is a pure
computation over an artifact, a public key, and an expected message. Pure
verification has no memory. If two isolated verifiers receive the same valid
artifact, each sees an unused artifact in its own local state. Both can accept.
No signature algorithm can tell one verifier what another verifier already did.

This paper studies that separation under a deliberately strong adversary. The
adversary controls the network and the orchestration service. It can choose the
actions shown to the signing interface and collect honest signatures on as many
other actions as it can persuade the human to approve. This is a
signature-soliciting adversary, not merely an attacker trying to forge a signature
without access to the signer. The construction should therefore prevent a real
signature for action $a'$ from being transplanted onto action $a$.

The contributions are:

\begin{enumerate}[leftmargin=*]
\item A small action-bound receipt construction and acceptance procedure that
separate portable evidence from stateful admission.
\item A security model in which an untrusted orchestrator can solicit honest
signatures over attacker-chosen actions.
\item An action-agreement argument under canonical encoding, collision resistance,
signature unforgeability, and uncompromised key pinning.
\item A simple impossibility result for global one-time admission by a static
artifact at isolated verifiers.
\item Two executable Tamarin models that verify the positive action-binding and
local-consumption properties and expose the two replay boundaries.
\end{enumerate}

The paper does not propose that a signed human interaction is itself an
authorization decision. A policy or authorization service decides whether the
evidence is sufficient. An executor then admits or refuses the exact action. The
receipt is an input to those decisions and later evidence of what was signed.

\section{System and Adversary Model}

\subsection{Parties}

The system has four roles.

\begin{description}[leftmargin=2.5cm,style=nextline]
\item[Approver $H$] A human associated with a signing key pair
$(sk_H,pk_H)$. The signing device requires a local user-verification event before
producing the signature.
\item[Orchestrator $O$] The service that prepares an action, requests approval,
and transports artifacts. It is not trusted for evidence integrity.
\item[Authorization service $A$] The policy decision point. It may accept a valid
receipt as one input to a grant. Its internal policy is outside the cryptographic
construction.
\item[Executor $X$] The component able to cause the consequential effect. It
checks the exact action and atomically records consumption before commitment.
\end{description}

The roles may be operated by the same organization, but the trust boundaries do
not disappear when software is colocated. In particular, a database row written
by $O$ is not treated as independent evidence of what $H$ signed.

\subsection{Signature-soliciting adversary}

The adversary controls all network traffic and the orchestrator. It may delay,
drop, reorder, duplicate, and replay messages. It chooses the action in any
signing request and receives every resulting receipt. It can therefore obtain
valid signatures over arbitrary actions that the human actually approves.

The adversary does not initially know $sk_H$. Key compromise is represented by an
explicit reveal event. Security statements are conditioned on no reveal before
the relevant acceptance. The relying party obtains the association between the
approver identifier and $pk_H$ through a trusted pinning process. The directory or
identity proof that establishes that association is not analyzed here.

This model is stronger than one that assumes the signer participates only in an
honest protocol run. The construction must remain action-bound even when the
attacker possesses honest signatures for many nearby actions.

\subsection{Action model}

An action $a$ is a typed object containing every field that can change the
security-relevant effect. For a payment this would include the asset, amount,
destination, source account, and relevant execution constraints. The action has
a deterministic canonical encoding $\Enc(a)$. JSON deployments can instantiate
this requirement with JCS~\cite{jcs}, provided the action schema also fixes
number, string, and omitted-field semantics.

\begin{definition}[Exact action]
Two actions are equal for authorization if and only if their canonical encodings
are byte-for-byte equal. An omitted material field is therefore a defect in the
action schema, not a fact repaired by the signature.
\end{definition}

The construction assumes that $\Enc$ is deterministic and injective over valid
actions. A parser differential or an ambiguous schema violates this assumption.
It must be addressed by canonicalization rules and conformance tests, not by the
symbolic signature proof.

\section{Construction}

\subsection{Signed authorization context}

For domain-separation tag $t$, action $a$, fresh nonce $n$, audience $aud$,
policy reference $p$, and validity interval $w$, define

\begin{equation}
  c = \Enc(t,\Hash(\Enc(a)),n,aud,p,w).
\end{equation}

The signing device renders the action from the same canonical object used to
compute the hash. After local user verification, it returns

\begin{equation}
  \sigma = \Sign_{sk_H}(c).
\end{equation}

The authorization receipt is

\begin{equation}
  R = (a,n,aud,p,w,H,\sigma,\mathit{keyref}).
\end{equation}

The full action travels with the receipt. This avoids asking a later verifier to
recover security-relevant parameters from operator prose. The key reference lets
the verifier resolve a pinned public key and its applicable status information.

\subsection{Offline evidence verification}

Given $R$, a relying party reconstructs $c$, resolves the pinned key, and checks:

\begin{enumerate}[leftmargin=*]
\item the action is well formed and canonically encoded;
\item the expected audience and policy reference equal the signed values;
\item the relevant validity and key-status inputs satisfy local policy; and
\item $\Verify_{pk_H}(c,\sigma)=1$.
\end{enumerate}

The signature and binding checks are offline if the verifier already has the
necessary key and status material. The word ``offline'' does not make those
inputs current. A verifier using an old revocation snapshot learns only what was
valid under that snapshot.

\subsection{Admission and consumption}

Verification produces evidence. Admission is a separate state transition. Let
$D$ be an authoritative consumption domain and let

\begin{equation}
  k_R = (H,n,aud)
\end{equation}

be the receipt's consumption key. The executor performs one atomic operation:

\begin{equation}
  \Accept_D(R,a_{exec}) =
  \begin{cases}
  \mathsf{refuse}, & a_{exec} \ne a,\\
  \mathsf{refuse}, & \text{verification or policy fails},\\
  \mathsf{refuse}, & k_R \text{ is already consumed in }D,\\
  \mathsf{commit}, & \Consume_D(k_R)\text{ succeeds atomically}.
  \end{cases}
\end{equation}

The consumption record is written only for the action that is about to be
committed. A refusal does not spend the authorization. The effect itself may need
an idempotency key or a transaction in the system of record. A consumption record
proves at-most-once admission by this gate; it does not prove an exactly-once
downstream physical or financial effect.

\subsection{Claim separation}

Table~\ref{tab:claims} states what each mechanism contributes. The separation is
important because a receipt can remain valid evidence after it has been consumed.
Consumption changes whether the action may be admitted again; it does not make the
historical signature false.

\begin{table}[h]
\centering
\small
\begin{tabular}{p{0.25\linewidth}p{0.29\linewidth}p{0.34\linewidth}}
\toprule
Mechanism & Establishes & Does not establish \\
\midrule
Signature over $c$ & Integrity and origin under the pinned key & Human identity, comprehension, policy sufficiency \\
Action hash and canonical encoding & Binding to the exact encoded action & Completeness of the action schema or honest display \\
Fresh nonce & Distinguishes authorization instances & Erasure or single use after disclosure \\
Audience and policy commitment & Scope of intended reliance & Current authority or policy correctness \\
Atomic consumption in $D$ & At-most-once admission inside $D$ & Global replay prevention outside $D$ or downstream effect \\
\bottomrule
\end{tabular}
\caption{Mechanisms and bounded claims.}
\label{tab:claims}
\end{table}

This also explains why verification and admission should have different outputs.
An offline verifier can return that the signature and action binding are valid
under supplied trust inputs. Only the authoritative executor can return that the
authorization is currently unused in its domain and has been admitted.

\begin{center}
\setlength{\tabcolsep}{6pt}
\begin{tabular}{>{\raggedright\arraybackslash}p{0.19\linewidth}c>{\raggedright\arraybackslash}p{0.24\linewidth}c>{\raggedright\arraybackslash}p{0.27\linewidth}}
\toprule
Orchestrator & & Approver device & & Authorization service and executor \\
\midrule
Prepare canonical $a$ & $\longrightarrow$ & Render $a$; require user verification & & \\
Receive $R$ & $\longleftarrow$ & Sign $c$ and return $R$ & & \\
Present $R,a$ & $\longrightarrow$ & & $\longrightarrow$ & Verify evidence; decide policy; atomically consume and commit \\
\bottomrule
\end{tabular}
\end{center}

\section{Security Properties}

\subsection{Action agreement}

\begin{definition}[Action agreement]
If a relying party accepts a receipt naming $H$ for action $a$, and $H$'s signing
key was not compromised before acceptance, then $H$ previously produced a
signature over a context containing $\Hash(\Enc(a))$ and the same nonce, audience,
policy reference, and validity interval.
\end{definition}

This definition is intentionally about the signed bytes. It does not claim that
the human understood the action, that the display was honest, or that policy
should authorize it.

\paragraph{Security experiment.}
The challenger generates $(sk_H,pk_H)$ and gives $pk_H$ to the adversary. The
adversary has a signing oracle that takes a valid action and context parameters,
records a user-verification event, and returns the corresponding receipt. This is
the cryptographic abstraction of an orchestrator that can solicit signatures on
chosen actions. The adversary wins if it outputs $(R,a)$ such that evidence
verification accepts under $pk_H$, but the exact context reconstructed from $a$
does not appear in the signing-oracle transcript. We call this experiment
$\mathsf{ABind}$.

The experiment is not intended to model whether the human should have approved
an oracle query. The oracle records that the approver did sign it. The security
goal is that a receipt from one query cannot be presented as evidence for a
different action or context.

\begin{proposition}[Cross-action non-transferability]
Assume that the signature scheme is existentially unforgeable under chosen
message attack, $\Hash$ is collision resistant, $\Enc$ is injective over valid
actions, the verifier reconstructs the signed context from the received action,
and the pinned public key belongs to the named approver. An adversary that obtains
honest signatures on actions of its choice cannot cause acceptance for a new
action $a$ unless the approver signed the corresponding context, the signing key
was compromised, or one of the assumptions fails.
\end{proposition}

\begin{proof}
Suppose a receipt is accepted for $a$ without a prior signature on its context.
The verifier accepted a signature under $pk_H$ for the reconstructed $c$. If the
signature was not previously returned on $c$, this is a signature forgery. If it
was returned for a context containing a different action $a'$, equality of the
verified contexts requires either $\Enc(a)=\Enc(a')$ or a hash collision. By
injectivity, the first case gives $a=a'$. The remaining cases contradict the
stated assumptions.
\end{proof}

The proposition gives a direct reduction from a successful $\mathsf{ABind}$
adversary to either signature forgery, a hash collision, or an encoding collision.
Domain separation in $t$ also prevents a signature produced for another message
type from being interpreted as an authorization receipt.

Fresh nonces stop one approval instance from being confused with another, but a
nonce does not erase a valid signature after use. That property needs state.

\subsection{Local injective admission}

\begin{definition}[Local injective admission]
Within a consumption domain $D$, no two distinct commitment events accept the
same receipt consumption key $k_R$.
\end{definition}

This follows from a linearizable compare-and-set, uniqueness constraint inside
the same transactional database, or equivalent atomic mechanism. Merely writing
a log entry after dispatch does not provide the property because two concurrent
requests may both pass the pre-check.

The ordering of checks matters. The executor first validates that the nonce is
available without spending it. It then validates the action, audience, policy,
and other required evidence. Only when it commits to the effect does it finalize
consumption. A refusal must not consume. A deployment may use a short reservation
while checks run, but every refusal and abandoned attempt must release that
reservation. Otherwise a rejected action can permanently destroy a valid grant.

\subsection{Why offline artifacts cannot be globally single use}

\begin{proposition}[No global single use at isolated verifiers]
Let $R$ be a static, valid authorization artifact addressed to an admission set
containing two verifiers $V_1$ and $V_2$. Assume the verifiers start from states
in which $R$ is locally unused, do not communicate before deciding, and each must
accept a fresh valid artifact in an otherwise valid request. No deterministic or
randomized local verification procedure can guarantee that at most one verifier
accepts $R$ while preserving acceptance at both verifiers when each is run alone.
\end{proposition}

\begin{proof}
Present $R$ to both verifiers. The local view of $V_1$ is indistinguishable from
an execution in which only $V_1$ receives the fresh artifact. Correctness
therefore requires $V_1$ to accept. The same argument applies to $V_2$. With
deterministic verifiers, both accept. With randomized verifiers that accept a
fresh valid artifact with nonzero probability, there is a nonzero probability
that both accept. Preventing the joint event requires information that changes
one local view, such as shared state, coordination, a unique online token
exchange, or an artifact whose audience permits only one admission domain.
\end{proof}

Audience restriction changes the set over which the claim is made; it does not
make offline state global. If several replicas share one audience and may each
cause the effect, they still need a common atomic consumption record or an
idempotent system-of-record operation.

\subsection{Available design choices}

Proposition 2 does not say that useful systems are impossible. It identifies the
resource each design must spend. Four common choices are:

\begin{enumerate}[leftmargin=*]
\item \textbf{Shared atomic state.} All executors consult one linearizable
consumption domain. This preserves flexible routing but adds an online dependency.
\item \textbf{Single-domain addressing.} The signed audience names one executor
or one domain with an internal atomic store. This preserves offline evidence
verification but gives up transparent admission at other domains.
\item \textbf{Idempotent effect key.} The system of record rejects a second
operation carrying the same key. This moves consumption to the effect boundary
and is often the strongest placement when the effect system supports it.
\item \textbf{Detect after use.} Independent domains accept offline and later
reconcile duplicate use. This provides double-use evidence, not at-most-once
prevention, and is inappropriate when the second effect is irreversible.
\end{enumerate}

Any claim of single use should therefore name both the consumption key and the
domain over which uniqueness is enforced. ``The receipt is non-replayable'' is
incomplete without those two facts.

\section{Symbolic Analysis}

\subsection{Method}

We use Tamarin Prover 1.10.0 with its standard symbolic signing and hashing
theories and Maude 3.4~\cite{tamarin}. The network is controlled by the
Dolev-Yao adversary~\cite{dolevyao}. The
public key is public. A relying party pins it out of band. The action in a signoff
request comes from an adversary-controlled input, so the model permits signatures
over arbitrary attacker-chosen actions.

User verification is represented by a linear fact consumed by the signing rule.
This models the requirement that the authenticator gate each signature on a local
verification event, as in a WebAuthn-style ceremony~\cite{webauthn}. It does not
model authenticator firmware, biometric quality,
or display fidelity. Canonical encoding is represented by injective symbolic
terms, and signatures have perfect symbolic cryptography.

The central signing rule emits the tuple containing the action, nonce, and
signature onto the adversarial network:

\begin{Verbatim}[fontsize=\small]
rule Human_Sign_Receipt:
  [ UVDone(H,a,n), !HumanLtk(H,ltk) ]
  --[ Signed(H,a,n) ]->
  [ Out(<a,n,sign(<'ep_signoff_v1',h(a),n>,ltk)>) ]
\end{Verbatim}

The action is not generated by an honest role. It originates in an \code{In(a)}
premise in the request rule, which means the adversary chooses it. This detail is
load-bearing: the action-agreement lemma is checked even when the adversary can
obtain honest signatures on other chosen actions.

\subsection{Core model}

The core model has two acceptance rules. The first checks the signature but does
not consume the receipt. The second emits a consumption event constrained to
occur once. Five lemmas are evaluated:

\begin{table}[h]
\centering
\small
\begin{tabular}{p{0.48\linewidth}p{0.19\linewidth}p{0.19\linewidth}}
\toprule
Lemma & Expected & Result \\
\midrule
Honest user-verified receipt is executable & witness & verified, 8 steps \\
Acceptance implies prior user verification and signature for the exact action & all traces & verified, 12 steps \\
Honest signatures for other actions cannot be transplanted & all traces & verified, 12 steps \\
Checked acceptance is injective with consumption & all traces & verified, 6 steps \\
Unchecked acceptance is injective & false & falsified, 10 steps \\
\bottomrule
\end{tabular}
\caption{Results for \code{ep\_receipt\_core.spthy}. The last lemma is retained
as a negative control.}
\end{table}

The falsified lemma has the expected trace: the adversary obtains one honest
receipt and presents it twice to the unchecked acceptance rule. The trace does not
forge a signature or change the action.

\subsection{Split-domain model}

The second model replaces the core model's single global consumption restriction
with a restriction indexed by domain $D$. Each domain is internally atomic, but
there is no fact or rule that coordinates different domains.

\begin{table}[h]
\centering
\small
\begin{tabular}{p{0.48\linewidth}p{0.19\linewidth}p{0.19\linewidth}}
\toprule
Lemma & Expected & Result \\
\midrule
Honest user-verified receipt is executable & witness & verified, 8 steps \\
Action binding survives local consumption & all traces & verified, 6 steps \\
Acceptance is injective inside one domain & all traces & verified, 2 steps \\
Acceptance is globally injective across domains & false & falsified, 8 steps \\
\bottomrule
\end{tabular}
\caption{Results for \code{ep\_receipt\_split\_domains.spthy}. The final
counterexample contains no key-reveal event.}
\end{table}

The counterexample is the symbolic form of Proposition 2. One honest signature is
delivered to two relying parties operating in different consumption domains.
Each domain consumes it once. Local injectivity holds, while global injectivity
fails. This is not a break of action binding. It is a failure to coordinate
state.

\subsection{Interpretation}

The models establish three points under their assumptions.

\begin{enumerate}[leftmargin=*]
\item Giving the adversary access to honest signatures over chosen actions does
not defeat exact action binding.
\item A signature check without consumption permits same-action replay.
\item Consumption state that is atomic only locally does not compose into a
global at-most-once property.
\end{enumerate}

The models do not prove computational security of a particular signature suite.
They also do not prove the directory, canonical parser, clock, authorization
policy, or downstream effect system. Those exclusions are part of the result,
because each excluded component supplies an assumption the theorem actually uses.

\section{Placement in an Agent Authorization Flow}

Action-bound approval is most useful when the need for human authorization is
discovered after an agent has begun work, a gap also identified in the current
AI-agent authorization Internet-Draft~\cite{klrc}. The agent pauses before the
consequential step. The orchestrator constructs the exact action object and
requests the receipt. An authorization service validates the receipt together
with delegation, identity, risk, and policy inputs. If it grants authority, the
executor still compares the action at commitment time and atomically consumes the
authorization.

This placement separates four questions that are easy to collapse:

\begin{enumerate}[leftmargin=*]
\item \textbf{Evidence validity:} does the signature verify over the exact action?
\item \textbf{Policy sufficiency:} does this receipt, with other evidence, satisfy
the relying party's rule?
\item \textbf{Authorization:} does the policy decision point grant this action?
\item \textbf{Admission:} may the executor commit it now, and has this grant
already been consumed in the relevant domain?
\end{enumerate}

The receipt answers the first question. It can be an input to the second and
third. It does not replace the fourth. A deployment that verifies the artifact
only in agent-side middleware can protect against agent error and network replay
through that path, but an operator who controls an alternate execution path can
bypass it. Strong at-most-once admission requires that every protected path meet
the authoritative consumer or system of record.

\section{Related Work}

Transaction authorization and what-you-see-is-what-you-sign systems bind a user
approval to transaction data rather than to a general authenticated
session~\cite{wysiwys}.
The present construction applies the same exact-message discipline to an agent
action discovered during execution, and places the operator that solicits the
signature inside the adversarial model.

OAuth Rich Authorization Requests~\cite{rar} and transaction
tokens~\cite{transactiontokens} carry structured authorization details and
workload identity across services. They are useful
inputs to the policy and delegation layers. They do not by themselves provide an
approver-held signature over the final action or a global consumption service.
OpenID Client-Initiated Backchannel Authentication~\cite{ciba} begins with a client request
and provides an out-of-band authentication and authorization flow. The problem
studied here is narrower: the evidence artifact and admission rule when the exact
need for human authorization appears during an agent run.

Formal analyses of OAuth 2.0~\cite{oauthformal} and FIDO2~\cite{fido2formal}
show the value of stating web and authenticator assumptions explicitly. This
paper treats WebAuthn-style user verification and key custody as assumptions and
verifies the receipt protocol around them. Tamarin has been used for symbolic
analysis of authentication, authorization, and payment protocols. Our models use
its trace logic to retain both the positive agreement theorem and the expected
replay counterexamples.

Recent work on consent integrity for black-box agents~\cite{consentintegrity}
focuses on the trusted path from a boundary action to the human display and back
to execution. That problem is complementary. The present paper assumes a fixed
action object, then asks what the resulting signature proves to a later verifier
and where single-use admission can be enforced.

Offline payment and electronic cash systems~\cite{chaum,offlinecash} address double spending with online
checks, tamper-resistant hardware, or later detection and identification. The
same fundamental distinction appears here in a simpler setting: a portable
cryptographic object can be copied. Preventing two isolated executors from acting
on the copy needs state or a restricted execution locus.

\section{Limitations}

The construction has several explicit limits.

\begin{itemize}[leftmargin=*]
\item \textbf{Human perception.} The signature covers canonical bytes, not human
understanding. A compromised renderer can mislead the approver unless the signing
surface is independently trusted or the action is displayed through a protected
path.
\item \textbf{Identity binding.} The symbolic proof pins a public key. It does not
prove that a directory correctly maps that key to a natural person or role.
\item \textbf{Key and status freshness.} Offline verification can use stale key
status. Current revocation generally requires a current authenticated input.
\item \textbf{Canonicalization.} The proof assumes an injective encoding. Parser
differentials, omitted material fields, and unsafe numeric representations are
implementation concerns that need separate testing.
\item \textbf{Policy.} A valid receipt may be insufficient under policy, or may
authorize a bad action. Cryptographic validity is not business correctness.
\item \textbf{Effects.} Atomic consumption provides at-most-once admission in its
domain. It does not prove that a remote effect occurred, occurred once, or can be
reversed.
\item \textbf{Availability and coercion.} The adversary can deny service, and the
construction does not detect coercion or rubber-stamping.
\end{itemize}

\section{Conclusion}

Action-bound signatures and one-time execution solve different problems. A
signature over a canonical action lets a third party verify, without trusting the
approval operator, which bytes an uncompromised approver key signed. That property
survives an adversary that can solicit honest signatures over other actions. It is
portable and can be checked offline.

Single use is not portable in the same way. A receipt can be copied, and isolated
verifiers cannot learn each other's decisions from a static artifact. At-most-once
admission therefore belongs at a shared atomic consumption boundary or an
equivalent uniquely addressed execution point. The Tamarin models preserve this
line: they verify action agreement and local injectivity, and they find the replay
traces when consumption is absent or split across domains.

For agent authorization, the practical rule is concise. Use the receipt as
evidence of exact human approval. Let the authorization service decide whether
that evidence is sufficient. Let the executor compare the exact action, consume
the grant atomically, and then commit. Evidence can travel. Admission needs state.

\section*{Artifact Availability}

The Tamarin models, exact tool versions, model hashes, and rerun commands are in
the public EMILIA Protocol repository under
\code{papers/eprint-action-bound-authorization}. The focused paper is a preprint
candidate and has not been submitted to the IACR Cryptology ePrint Archive as of
the date on its title page.

\begin{thebibliography}{99}

\bibitem{dolevyao}
D. Dolev and A. C. Yao.
On the security of public key protocols.
\emph{IEEE Transactions on Information Theory}, 29(2), 1983.

\bibitem{tamarin}
S. Meier, B. Schmidt, C. Cremers, and D. Basin.
The Tamarin prover for the symbolic analysis of security protocols.
In \emph{CAV 2013}, LNCS 8044, 2013.

\bibitem{oauthformal}
D. Fett, R. K\"usters, and G. Schmitz.
A comprehensive formal security analysis of OAuth 2.0.
In \emph{ACM CCS 2016}, 2016.

\bibitem{fido2formal}
M. Barbosa, A. Boldyreva, S. Chen, and B. Warinschi.
Provable security analysis of FIDO2.
In \emph{CRYPTO 2021}, LNCS 12827, pages 125--156, 2021.

\bibitem{wysiwys}
P. Landrock and T. Pedersen.
WYSIWYS? What you see is what you sign?
\emph{Information Security Technical Report}, 3(2):55--61, 1998.

\bibitem{jcs}
A. Rundgren, B. Jordan, and S. Erdtman.
JSON Canonicalization Scheme (JCS).
RFC 8785, June 2020.

\bibitem{webauthn}
W3C.
Web Authentication: An API for accessing public key credentials, Level 2.
W3C Recommendation, April 2021.

\bibitem{rar}
T. Lodderstedt, J. Richer, and B. Campbell.
OAuth 2.0 Rich Authorization Requests.
RFC 9396, May 2023.

\bibitem{ciba}
OpenID Foundation.
OpenID Connect Client-Initiated Backchannel Authentication Flow, Core 1.0.
September 2021.

\bibitem{transactiontokens}
J. Richer, P. Hunt, and M. B. Jones.
Transaction Tokens.
Internet-Draft, work in progress.

\bibitem{klrc}
P. Kasselman, J. Lombardo, Y. Rosomakho, B. Campbell, N. Steele, and A. Parecki.
AI Agent Authentication and Authorization.
Internet-Draft \code{draft-klrc-aiagent-auth-03}, work in progress, July 2026.

\bibitem{consentintegrity}
X. Weng.
What You Approve Is What Executes: Consent Integrity for Black-Box LLM Agents.
arXiv:2606.02668, June 2026.

\bibitem{chaum}
D. Chaum.
Blind signatures for untraceable payments.
In \emph{CRYPTO 1982}, 1983.

\bibitem{offlinecash}
S. Brands.
\emph{Rethinking Public Key Infrastructures and Digital Certificates: Building
in Privacy}.
MIT Press, 2000.

\bibitem{emiliareceipt}
I. Schrock.
Authorization Receipts for High-Risk Agent Actions.
Internet-Draft \code{draft-schrock-ep-authorization-receipts}, work in progress,
August 2026.

\end{thebibliography}

\end{document}
